% Generated draft: requires mathematical and source review.
\documentclass[11pt]{article}
\usepackage[margin=25mm]{geometry}
\usepackage{fontspec}
\setmainfont{Latin Modern Roman}
\usepackage{amsmath,amssymb,mathtools,unicode-math}
\setmathfont{Latin Modern Math}
\usepackage{enumitem,needspace,xurl,hyperref}
\setlength{\parindent}{0pt}
\setlength{\parskip}{5pt}
\setlength{\emergencystretch}{3em}
\newcommand{\N}{\mathbb{N}}
\newcommand{\Z}{\mathbb{Z}}
\newcommand{\Q}{\mathbb{Q}}
\newcommand{\R}{\mathbb{R}}
\newcommand{\C}{\mathbb{C}}
\newcommand{\F}{\mathbb{F}}
\newcommand{\A}{\mathbb{A}}
\newcommand{\PP}{\mathbb{P}}
\newcommand{\E}{\mathbb{E}}
\newcommand{\eps}{\varepsilon}
\newcommand{\dd}{\,\mathrm{d}}
\begin{document}
\section*{UnsolvedMath ID: 800660}
\textbf{Source problem code:} catalogue entry
\subsection{Definitions and mathematical statement}
The supplied record asks about two selection principles involving \(\mathcal D_\Omega\), but does not define this family or specify the ambient objects to which the properties apply. Consequently, the following gives the standard abstract selection and game formulations; identifying their intended parameters remains necessary for a complete mathematical statement.

Let \(U\) be a set and let \(\mathscr A\subseteq\mathcal P(U)\setminus\{\varnothing\}\). These are auxiliary parameters, not an asserted definition of \(\mathcal D_\Omega\). Define
\[
\mathsf S_1(\mathscr A,\mathscr A)
\quad\Longleftrightarrow\quad
\forall (A_n)_{n\in\N}\in\mathscr A^{\N}\ 
\exists (a_n)_{n\in\N}\in U^{\N}\ 
\left[(\forall n\in\N\;a_n\in A_n)\ \land\ \{a_n:n\in\N\}\in\mathscr A\right].
\]
Define
\[
\mathsf S_{\mathrm{fin}}(\mathscr A,\mathscr A)
\quad\Longleftrightarrow\quad
\forall (A_n)_{n\in\N}\in\mathscr A^{\N}\ 
\exists (F_n)_{n\in\N}\ 
\left[(\forall n\in\N\;F_n\subseteq A_n\text{ is finite})\ \land\ \bigcup_{n\in\N}F_n\in\mathscr A\right].
\]
Here \(\N=\{0,1,2,\ldots\}\), and finite selections are allowed to be empty in this template.

In the corresponding game \(\mathsf G_1(\mathscr A,\mathscr A)\), inning \(n\) consists of ONE choosing \(A_n\in\mathscr A\), followed by TWO choosing \(a_n\in A_n\). TWO wins exactly when \(\{a_n:n\in\N\}\in\mathscr A\); otherwise ONE wins. In \(\mathsf G_{\mathrm{fin}}(\mathscr A,\mathscr A)\), TWO instead chooses a finite subset \(F_n\subseteq A_n\), and wins exactly when \(\bigcup_{n\in\N}F_n\in\mathscr A\). A strategy for ONE specifies a legal move at each finite legal history at which ONE moves, and is winning if every infinite legal play following it is won by ONE.

With \(\mathscr A\) replaced by the intended, presently unspecified \(\mathcal D_\Omega\), the two questions are whether
\[
\mathsf S_1(\mathcal D_\Omega,\mathcal D_\Omega)
\quad\overset{?}{\Longleftrightarrow}\quad
\text{ONE has no winning strategy in }\mathsf G_1(\mathcal D_\Omega,\mathcal D_\Omega)
\]
and whether
\[
\mathsf S_{\mathrm{fin}}(\mathcal D_\Omega,\mathcal D_\Omega)
\quad\overset{?}{\Longleftrightarrow}\quad
\text{ONE has no winning strategy in }\mathsf G_{\mathrm{fin}}(\mathcal D_\Omega,\mathcal D_\Omega).
\]
These are separate proposed equivalences. Their intended scope cannot be fixed from the supplied record: the definition of \(\mathcal D_\Omega\), the class of ambient objects, and any additional hypotheses are missing. The game conventions above are the standard abstract template and require confirmation against the cited source; they are not recovered source definitions.
\subsection{Short English statement}
For each of two selection properties involving \(\mathcal D_\Omega\), ask whether the property is equivalent to ONE having no winning strategy in its corresponding infinite selection game. The supplied record omits the definition of \(\mathcal D_\Omega\) and the ambient hypotheses, so this draft gives a conditional formal template requiring human review.
\subsection{Sources}
\begin{itemize}
\item UnsolvedMath local registry, permanent ID \(800660\), exact source problem code \texttt{catalogue entry}; supplied title: ``Selection Principles Milestone Problems --- Problem 33 ([coc5, Problem 3]).'' The registry attributes the record to Tsaban, with proposed year 2003, and labels it open. This status is supplied metadata, not an independently verified mathematical or literature assessment.
\item Supplied source list: Tsaban, \emph{Selection Principles in Mathematics: A Milestone of Open Problems} (2003), Problem 33. Supplied URL: \url{https://arxiv.org/abs/math/0312182}. The record reports access on 2026-10-16 and extraction by \texttt{arxiv\_latex}; these are supplied provenance details, not a claim of independent access or verification.
\item Supplied attribution/context: [coc5, Problem 3]. The bibliographic identity of \texttt{coc5} is not included in the supplied record, so the underlying reference and its precise hypotheses remain unresolved.
\end{itemize}
\end{document}
